% This document was generated by an LLM.

\documentclass{essay}

\newcommand{\Irr}{\R \setminus \Q}
\DeclareMathOperator{\diam}{diam}

\title{Four Sizes of the Rational Numbers}
\author{}
\date{}

\begin{document}

\maketitle

\section{The question is not one question}

The rational numbers are dense in the real line, and so are the
irrationals. The rationals are countable, and the irrationals are not.
These facts are usually assembled into a slogan: there are
\emph{infinitely more} irrationals than rationals. The slogan is not
merely loose. It is not a statement, because ``infinitely more'' names
no operation on cardinals, and, as will appear in Section~\ref{sec:card},
the question it gestures at has no answer inside ZFC.

The purpose of this essay is to replace the slogan with something
better. The claim to be defended is that \emph{size} of a subset of
$\R$ is not a single notion but at least four inequivalent ones ---
density, cardinality, measure, and category --- and that $\Q$ is the
example on which they visibly come apart. Under three of them $\Q$ is
as small as a nonempty set can be. Under the fourth it is as large as
any subset of $\R$ can be. No single intuition survives this.

\section{Setting the two density claims straight}

The density statements are frequently quoted in a weakened form:
between any two rationals there is an irrational, and between any two
irrationals there is a rational. Both are true, but neither is the
property being named. Density of a set $S \subseteq \R$ means that
every nonempty open interval of $\R$ meets $S$ --- the endpoints range
over all of $\R$, not over a distinguished subset. The restricted
versions happen to be equivalent to the real ones, but only by a
detour through the real ones, which makes them a poor way to state the
theorem.

\begin{theorem}[Density of $\Q$]
For all real $a < b$ there is $q \in \Q$ with $a < q < b$.
\end{theorem}

\begin{proof}
By the Archimedean property choose $n \in \N$ with $n(b - a) > 1$, so
that $1/n < b - a$. Let $m = \lfloor na \rfloor + 1$, so that
$m - 1 \leq na < m$. Then $a < m/n$, and
\[
  \frac{m}{n} \leq \frac{na + 1}{n} = a + \frac{1}{n} < a + (b - a) = b .
\]
Hence $q = m/n$ works.
\end{proof}

\begin{theorem}[Density of $\Irr$]
For all real $a < b$ there is $x \in \Irr$ with $a < x < b$.
\end{theorem}

\begin{proof}
Apply the previous theorem to the interval $(a - \sqrt{2},\, b -
\sqrt{2})$, obtaining $q \in \Q$ with $a - \sqrt{2} < q < b -
\sqrt{2}$. Then $q + \sqrt{2} \in (a,b)$, and it is irrational: were
$q + \sqrt{2} = r \in \Q$, we would have $\sqrt{2} = r - q \in \Q$.
\end{proof}

The second proof is worth a moment. It uses nothing about $\Q$ except
that it is a dense subgroup of $(\R,+)$ and that its complement is
nonempty; the translate of a dense set by a fixed real is dense, and
the translation carries $\Q$ off itself. The same argument shows that
any dense set has dense complement whenever the complement contains a
coset of the set. This is the first sign that density is a coarse
instrument.

\section{Density carries no size information at all}

\begin{proposition}
Density is independent of cardinality: each of the four combinations of
(countable, uncountable) with (dense, nowhere dense) is realised by a
subset of $\R$.
\end{proposition}

\begin{proof}[Discussion]
$\Q$ is countable and dense. $\Irr$ is uncountable and dense. The set
$\Z$ is countable and nowhere dense. The Cantor middle-thirds set $C$
is uncountable --- it surjects onto $\{0,1\}^{\N}$ by ternary
expansion --- and nowhere dense, being closed with empty interior.
\end{proof}

So the observation that $\Q$ and $\Irr$ are both dense tells us
precisely nothing about their relative size, and one should resist the
common instinct to treat ``dense'' as a largeness property at all.
Density is a statement about closure: $S$ is dense exactly when
$\overline{S} = \R$. It says that $S$ approximates every real
arbitrarily well. It does not say that $S$ has many elements, and the
Cantor set shows the converse failure --- a set can have the
cardinality of the continuum while approximating almost nothing.

\section{Cardinality is an equivalence class, not a magnitude}
\label{sec:card}

\begin{theorem}
$\Q$ is countable and $\Irr$ is not.
\end{theorem}

\begin{proof}
Each nonzero $q \in \Q$ has a unique representation $\pm m/n$ in lowest
terms with $m, n \in \N$, which gives an injection $\Q \to \Z \times
\N$; the latter is countable, hence so is $\Q$. For the second claim,
suppose $\Irr$ were countable. Since $\R = \Q \cup (\Irr)$ is then a
union of two countable sets, $\R$ would be countable, contradicting
Cantor's theorem that no surjection $\N \to \R$ exists.
\end{proof}

\begin{corollary}
$|\Irr| = |\R| = 2^{\aleph_0}$.
\end{corollary}

Now the point of the section. The relation $\aleph_0 < 2^{\aleph_0}$
unpacks to exactly two assertions: there is an injection $\Q \to \Irr$,
and there is no bijection. Nothing further is being claimed, and in
particular no quantity has been produced that measures the gap. The
temptation to ask ``how much bigger'' has two independent obstructions.

The first is algebraic. Cardinal addition absorbs: $\aleph_0 +
2^{\aleph_0} = 2^{\aleph_0}$, which is just the statement that
adjoining $\Q$ back to $\Irr$ recovers $\R$ without changing its
cardinality. Because addition absorbs, it is not cancellative, and
there is no subtraction. The expression ``how many more irrationals
than rationals'' is therefore not an expression.

The second obstruction is deeper. Even the question of \emph{where}
$2^{\aleph_0}$ sits in the hierarchy of alephs is not settled by the
axioms. The continuum hypothesis, $2^{\aleph_0} = \aleph_1$, was shown
consistent with ZFC by G\"odel (1940) and its negation shown consistent
by Cohen (1963). So ``how much larger is the continuum than
$\aleph_0$'' has no answer in ZFC --- not an unknown answer, but no
answer determined by the theory. A phrase like ``infinitely more'' does
not merely overstate what has been proved; it presupposes a
determinate quantity that the axioms decline to fix.

What cardinality does give is a genuine and useful dichotomy:
countable sets can be enumerated, indexed, and exhausted by induction;
uncountable ones cannot. That is a structural fact with real
consequences, and it is all that $\aleph_0 < 2^{\aleph_0}$ asserts.

\section{Measure supplies what cardinality withholds}

Cardinality is invariant under arbitrary bijections and so cannot see
the geometry of the line. Lebesgue measure can.

\begin{theorem}
$\Q$ has Lebesgue measure zero.
\end{theorem}

\begin{proof}
Enumerate $\Q = \{q_1, q_2, \dots\}$ and fix $\varepsilon > 0$. Put
$I_k = (q_k - \varepsilon 2^{-k-1},\, q_k + \varepsilon 2^{-k-1})$.
Then $\Q \subseteq \bigcup_k I_k$ and
\[
  \sum_{k=1}^{\infty} |I_k| = \sum_{k=1}^{\infty} \varepsilon 2^{-k}
  = \varepsilon .
\]
Since $\varepsilon$ was arbitrary, the outer measure of $\Q$ is $0$.
\end{proof}

\begin{corollary}
$\lambda\big((\Irr) \cap [0,1]\big) = 1$.
\end{corollary}

This is the theorem behind the slogan ``almost every real is
irrational,'' and unlike the slogan about cardinality it has exact
content: a point chosen uniformly at random from $[0,1]$ lies in $\Q$
with probability zero. Note also what the covering argument used ---
only that $\Q$ is countable. Every countable set is null. The
implication is one-directional, and this matters:

\begin{example}
The Cantor set $C$ is null. Its complement in $[0,1]$ has measure
$\sum_{n \geq 0} 2^n 3^{-n-1} = 1$. Yet $|C| = 2^{\aleph_0}$.
\end{example}

So measure and cardinality separate in both directions: $\Q$ is small
in both senses, $C$ is null but of full cardinality, and $\Irr$ is
large in both. Three of the four combinations already occur among sets
we have named. A ``size'' that assigns $\Q$ and $C$ the same value
while cardinality distinguishes them, and distinguishes $C$ from
$[0,1]$ while cardinality does not, is plainly measuring something
else.

\section{Category gives a third answer, and it is incompatible}

\begin{definition}
A set is \emph{meagre} (first category) if it is a countable union of
nowhere dense sets, and \emph{comeagre} if its complement is meagre.
\end{definition}

\begin{proposition}
$\Q$ is meagre and $\Irr$ is comeagre.
\end{proposition}

\begin{proof}
$\Q = \bigcup_{q \in \Q} \{q\}$, a countable union of singletons, each
nowhere dense.
\end{proof}

By the Baire category theorem $\R$ is not meagre, so the notion is not
vacuous: meagre is a genuine smallness property, and it again places
$\Q$ on the small side. Three notions, three verdicts, all agreeing
that $\Q$ is negligible. It would be natural to conclude that measure
and category are two routes to the same underlying idea. They are not.

\begin{theorem}[Erd\H{o}s--Sierpi\'nski duality, decomposition form]
There is a partition $\R = A \sqcup B$ with $A$ meagre and $\lambda(B)
= 0$.
\end{theorem}

\begin{proof}[Sketch]
Enumerate $\Q = \{q_k\}$ and for each $n$ let $U_n = \bigcup_k (q_k -
2^{-k-n},\, q_k + 2^{-k-n})$, an open dense set of measure at most
$2^{1-n}$. Put $B = \bigcap_n U_n$. Then $\lambda(B) = 0$, while
$A = \R \setminus B = \bigcup_n (\R \setminus U_n)$ is a countable
union of closed sets with empty interior, hence meagre.
\end{proof}

The consequence is worth stating flatly. Every real number lies either
in a set that topology calls negligible or in a set that measure calls
negligible. The two notions of smallness are not refinements of one
another; they are, on the line, complementary. Any argument of the form
``this set is small, therefore its complement is large'' must therefore
name which smallness it means, and cannot silently transfer the
conclusion to the other.

\section{A fourth notion, for contrast}

Hausdorff dimension refines measure by asking how the covering cost
scales. For $s \geq 0$ one sets
\[
  \mathcal{H}^{s}(E) = \lim_{\delta \to 0}\ \inf\Big\{ \sum_i
  (\diam U_i)^{s} \ :\ E \subseteq \bigcup_i U_i,\ \diam U_i <
  \delta \Big\},
\]
and $\dim_H E$ is the critical exponent at which $\mathcal{H}^s(E)$
drops from $\infty$ to $0$. Every countable set has dimension $0$, so
$\dim_H \Q = 0$. But $\dim_H C = \log 2 / \log 3 \approx 0.631$.
Lebesgue measure assigns $\Q$ and $C$ the same value; dimension
separates them. Where cardinality was too coarse to see geometry,
measure is too coarse to see scaling.

\section{What each notion is for}

The four notions are not competitors, and none is the ``real'' one.
They answer different questions, and the right one is fixed by the use.

\begin{itemize}
\item \emph{Density} asks whether the set approximates every point. It
is the right notion when the object of interest is a continuous
function, since a continuous function is determined by its values on a
dense set --- which is exactly why $\Q$ being countable and dense makes
$\R$ separable and makes numerical analysis possible.

\item \emph{Cardinality} asks whether the set can be enumerated. It is
the right notion when the argument proceeds by listing, indexing, or
induction along an enumeration --- as in the covering proof above,
which needed countability and nothing else.

\item \emph{Measure} asks what the probability of landing in the set
is. It is the right notion for anything probabilistic or
integral-theoretic, and it is what licenses ``almost every.''

\item \emph{Category} asks whether the set is topologically typical. It
is the right notion for genericity arguments --- the standard proof
that the generic continuous function on $[0,1]$ is nowhere
differentiable is a category argument, and no measure argument
substitutes for it.
\end{itemize}

\section{Conclusion}

The four opening claims are true, but the moral usually drawn from them
is not. It is not that the irrationals outnumber the rationals by some
infinite factor; that statement has no referent, and the quantity it
implies is independent of the axioms. The moral is that $\Q$ is a set
on which our notions of size visibly disagree: dense, hence
topologically ubiquitous; countable, hence enumerable; null, hence
probabilistically invisible; meagre, hence topologically negligible ---
and, by the duality theorem, negligibility in the last two senses is
not one property but two opposed ones.

The instinct being corrected is the instinct to hear ``smaller'' as a
single word. On subsets of the real line it is at least four words, and
the rationals are where the difference is impossible to overlook.

\end{document}
